\begin{statementbox}
\textbf{条件与记号。} 椭圆（圆的情形也适用）及直线为
\[
 \frac{x^2}{\alpha}+\frac{y^2}{\beta}=1,\quad \alpha,\beta>0;
 \qquad Ax+By+C=0,\quad A^2+B^2>0.
\]
其中 $\alpha,\beta$ 是\textbf{分母，即半轴长的平方}；不预设谁大。
定义 $D=A^2\alpha+B^2\beta>0$，$H=D-C^2$。
交点记为 $P(x_1,y_1),Q(x_2,y_2)$，不与系数 $A,B$ 混用。

\textbf{消去 $y$：要求 $B\ne0$。}
\[
 Dx^2+2AC\alpha x+\alpha(C^2-B^2\beta)=0.
\]
\[
 \boxed{x_1+x_2=-\frac{2AC\alpha}{D}},\qquad
 \boxed{x_1x_2=\frac{\alpha(C^2-B^2\beta)}{D}}.
\]
\textbf{消去 $x$：要求 $A\ne0$。}
\[
 Dy^2+2BC\beta y+\beta(C^2-A^2\alpha)=0.
\]
\[
 \boxed{y_1+y_2=-\frac{2BC\beta}{D}},\qquad
 \boxed{y_1y_2=\frac{\beta(C^2-A^2\alpha)}{D}}.
\]
以上实交点坐标关系要求 $H\ge0$；相切时按 $P=Q$ 重计切点。

\textbf{位置关系：} $H>0$ 为两个不同交点，$H=0$ 为相切，$H<0$ 为无公共点。
在各自\textbf{消元条件满足时}，对应的判别式为
\[
 \Delta_x=4\alpha\beta B^2H\quad(B\ne0),\qquad
 \Delta_y=4\alpha\beta A^2H\quad(A\ne0).
\]
\textbf{中点与弦长：} 当 $H>0$ 时，水平、竖直及斜弦均有
\[
 M=\left(-\frac{AC\alpha}{D},-\frac{BC\beta}{D}\right),\qquad
 PQ=\frac{2\sqrt{\alpha\beta(A^2+B^2)H}}{D}.
\]
$H=0$ 时 $M$ 为切点，上式长度为 $0$，不再是一条非零弦。
若采用斜率公式 $PQ=\sqrt{1+k^2}\,|x_2-x_1|$，则必须有 $B\ne0$，$k=-A/B$。
\end{statementbox}
